\documentclass[11pt,a4paper]{article}
% =====================================================================
%  Shared preamble for the Part V lecture notes (lectures 19-22).
%
%  These four lectures sit outside Geron, so the notes ARE the primary
%  source and are examined on the same terms as the textbook parts.
%  Everything here is chosen to make that readable on paper: one column,
%  generous measure, and the same six-colour palette as the slides so a
%  student moving between the two is not re-learning what red means.
%
%  Build with pdflatex (twice, for the toc and the refs):
%      cd notes && latexmk -pdf lecture-19.tex
%  or  make -C notes
% =====================================================================
\usepackage[T1]{fontenc}
\usepackage[utf8]{inputenc}
\usepackage{newtxtext,newtxmath}      % Times-like: prints well, reads well
\usepackage[scaled=0.86]{inconsolata} % code, at a size that sits beside Times
\usepackage{microtype}
% newtx defines \Bbbk, \openbox, \proof and \endproof; amssymb and amsthm
% then redefine them and abort. Freeing the four names before those two
% packages load is the standard fix, and it costs nothing here: we use
% amsthm only for \newtheorem.
\let\Bbbk\relax
\usepackage{amsmath,amssymb}
\let\openbox\relax
\let\proof\relax
\let\endproof\relax
\usepackage{amsthm}
\usepackage{booktabs}
\usepackage{enumitem}
\usepackage{xcolor}
\usepackage{tcolorbox}
\usepackage{titlesec}
\usepackage{graphicx}
\usepackage[hidelinks]{hyperref}
\tcbuselibrary{skins,breakable}

% ---- the course palette, from assets/css/custom.css -------------------
\definecolor{aimlprimary}{HTML}{0B3D62}   % deep blue  - structure
\definecolor{aimlaccent} {HTML}{C0392B}   % brick red  - failures
\definecolor{aimlsuccess}{HTML}{14663A}   % green      - repairs
\definecolor{aimlmath}   {HTML}{6C3483}   % purple     - the derivation
\definecolor{aimlmuted}  {HTML}{4B5563}
\definecolor{aimlrule}   {HTML}{B0BCC7}
\definecolor{aimlsoft}   {HTML}{F4F7F9}

\usepackage[a4paper,margin=32mm,top=30mm,bottom=30mm]{geometry}
\setlength{\parskip}{0.45\baselineskip}
\setlength{\parindent}{0pt}
\linespread{1.06}

% ---- headings --------------------------------------------------------
\titleformat{\section}{\Large\bfseries\color{aimlprimary}}{\thesection}{0.7em}{}
\titleformat{\subsection}{\large\bfseries\color{aimlprimary}}{\thesubsection}{0.7em}{}
\titleformat{\subsubsection}{\bfseries\color{aimlmuted}}{}{0em}{}
\titlespacing*{\section}{0pt}{1.6\baselineskip}{0.5\baselineskip}
\titlespacing*{\subsection}{0pt}{1.2\baselineskip}{0.35\baselineskip}

% ---- boxes -----------------------------------------------------------
% One box per role, and the role is the colour. A student who has seen the
% slides already knows what each one means before reading a word of it.
\newtcolorbox{keybox}[1][]{
  breakable, enhanced, colback=aimlsoft, colframe=aimlprimary,
  boxrule=0.4pt, left=8pt, right=8pt, top=6pt, bottom=6pt,
  fonttitle=\bfseries\small, coltitle=white, colbacktitle=aimlprimary,
  attach boxed title to top left={yshift=-2mm, xshift=4mm},
  boxed title style={boxrule=0pt, arc=1pt}, #1}

\newtcolorbox{failbox}[1][]{
  breakable, enhanced, colback=white, colframe=aimlaccent,
  boxrule=0.9pt, left=8pt, right=8pt, top=6pt, bottom=6pt,
  fonttitle=\bfseries\small, coltitle=white, colbacktitle=aimlaccent,
  attach boxed title to top left={yshift=-2mm, xshift=4mm},
  boxed title style={boxrule=0pt, arc=1pt}, #1}

\newtcolorbox{derivbox}[1][]{
  breakable, enhanced, colback=white, colframe=aimlmath,
  boxrule=0.6pt, left=8pt, right=8pt, top=6pt, bottom=6pt,
  fonttitle=\bfseries\small, coltitle=white, colbacktitle=aimlmath,
  attach boxed title to top left={yshift=-2mm, xshift=4mm},
  boxed title style={boxrule=0pt, arc=1pt}, #1}

% An exercise. Deliberately the same shape as the notebook's `try` field:
% one change, and what should happen to the output.
\newtcolorbox{trybox}[1][]{
  breakable, enhanced, colback=white, colframe=aimlsuccess,
  boxrule=0.5pt, left=8pt, right=8pt, top=6pt, bottom=6pt,
  fonttitle=\bfseries\small, coltitle=white, colbacktitle=aimlsuccess,
  attach boxed title to top left={yshift=-2mm, xshift=4mm},
  boxed title style={boxrule=0pt, arc=1pt}, #1}

% ---- small semantics -------------------------------------------------
\newcommand{\scope}[1]{\textcolor{aimlmuted}{\small #1}}
\newcommand{\term}[1]{\textbf{#1}}
\newcommand{\code}[1]{\texttt{\small #1}}
\newcommand{\lecture}[1]{Lecture~#1}

\newtheorem{definition}{Definition}[section]
\newtheorem{proposition}[definition]{Proposition}

% ---- title block -----------------------------------------------------
% \notestitle{19}{Information retrieval: the lexical foundation}{SciFact (BEIR)}
\newcommand{\notestitle}[3]{%
  \thispagestyle{empty}
  \begin{flushleft}
    {\color{aimlmuted}\small\sffamily
     APPLICATIONS OF MACHINE LEARNING \quad$\cdot$\quad
     BSc Mathematics of Artificial Intelligence}\\[2mm]
    {\color{aimlrule}\rule{\textwidth}{0.8pt}}\\[4mm]
    {\color{aimlmuted}\large Lecture #1 \quad$\cdot$\quad Part V \quad$\cdot$\quad
     Extended notes}\\[3mm]
    {\Huge\bfseries\color{aimlprimary}\baselineskip=1.15\baselineskip #2\par}\vspace{4mm}
    {\color{aimlmuted}\large Dataset: #3}\\[3mm]
    {\color{aimlrule}\rule{\textwidth}{0.8pt}}
  \end{flushleft}
  \vspace{2mm}
  \begin{keybox}[title={Why these notes exist}]
    Lectures 19--22 sit outside G\'eron, the course's primary text. For those
    four lectures \emph{these notes are the primary source}, and they are
    examined on exactly the same terms as the textbook parts. Everything in
    the slide deck for this lecture is here, with the arguments written out at
    length rather than compressed onto a slide; the numbers are the ones the
    accompanying notebook prints, and every one of them is a measurement on
    the corpus named above rather than a figure quoted from a paper.
  \end{keybox}
  \vspace{1mm}
  {\small\tableofcontents}
  \vspace{2mm}
  \noindent{\color{aimlrule}\rule{\textwidth}{0.4pt}}
  \vspace{1mm}
}


\begin{document}
\notestitlefor{13}{III}{Transfer learning}{Flowers102}{Chapter 12}

\section{The condition, named}

\begin{center}
\small
\begin{tabular}{lr}
\toprule
\textbf{Model} & \textbf{Test accuracy} \\
\midrule
Always the commonest species & 3.87\% \\
Convolutional network from scratch, 4{,}807{,}494 parameters & 15.3\% \\
What the operator needs & 90\% \\
\bottomrule
\end{tabular}
\end{center}

Lecture 12 built that network and it is not close --- and its architecture is
not the problem, since the same shape of network solves Fashion-MNIST
comfortably.

\begin{keybox}[title={The diagnosis, in one sentence}]
We asked 1{,}020 labelled images to pay for 4{,}807{,}494 parameters, most of
which encode facts about \emph{photographs in general} --- edges, textures,
colour gradients --- rather than facts about flowers.

The repair is therefore not a better architecture, a longer run or a cleverer
schedule. It is starting from weights that already work.
\end{keybox}

Start with what is \emph{not} wrong. The training loop is the same five lines
that worked before. The metric is accuracy, anchored against two baselines. The
split came with the dataset and the test set was touched once. The architecture
is the textbook's own shape. Nothing here is a bug --- which is what makes this
a different kind of repair from Lecture 2's.

The measurement that names the problem is the gap: 71.8\% on the training set
against 13.7\% on the validation set, 58 points apart. Lecture 5 gave that
shape a name --- a persistent gap between two plateaus is \emph{variance}, not
bias. The network fits the training set and is still climbing; there is nothing
wrong with its capacity.

Lecture 5 gave two cures for variance. More training data: no, labelling is the
whole cost. More constraint on the hypothesis space: yes, and we have done it
once already --- weight sharing was a constraint and it bought a factor of
5{,}478{,}275. The question is where the next constraint comes from.

\subsection{What a pretrained backbone has actually learned}

Not ``flowers''. A hierarchy, and only the bottom of it is general.

\begin{center}
\small
\begin{tabular}{lll}
\toprule
\textbf{Depth} & \textbf{What the filters respond to} & \textbf{Transfers?} \\
\midrule
first block & edges, colour opponency, simple gradients & to essentially
                                                          anything \\
middle blocks & textures, repeated motifs, part-like shapes & to most natural
                                                              images \\
last blocks & object-specific configurations & only to similar objects \\
the head & the 1{,}000 ImageNet classes & not at all --- it is replaced \\
\bottomrule
\end{tabular}
\end{center}

Which is the whole recipe: keep the bottom, adapt the top, replace the head.

\begin{keybox}[title={Freezing is not a compromise forced by compute}]
It is a statement that those layers already encode what you would otherwise be
trying to learn from 1{,}020 images.
\end{keybox}

Lecture 12 ended with a first layer whose filters were still noise after 364
seconds. The same layer from a network trained on ImageNet shows oriented
light--dark boundaries at every angle and opponent-colour blobs --- 9{,}408
numbers, and not one of them is about flowers, or about any of the 1{,}000
classes that network was fitted on. They are about \emph{photographs}.

ImageNet is 1.28 million labelled photographs over 1{,}000 classes, none of
them our species. Somebody already spent that on learning edges, colours,
textures and parts, and those weights are a download. Our 1{,}020 labels should
be spent on the only thing specific to us: which combination of parts is which
species.

\section{Freeze, unfreeze, and augment}

\subsection{Probe first, fine-tune second}

Two decisions, in this order. A \term{linear probe} freezes everything,
extracts features once and fits a linear classifier --- minutes, and it tells
you whether the backbone understands your images at all. Then \term{fine-tune}:
unfreeze the last block or two at a small learning rate, and only if the probe
was promising.

\begin{keybox}[title={Why the order is not arbitrary}]
The probe is a measurement of the \emph{features}, uncontaminated by anything
you do to them. If it barely beats the trivial baseline, the features are wrong
for this domain and fine-tuning is polishing the wrong object. If it is already
close, fine-tuning has little left to buy and may cost you generalisation.
\end{keybox}

How much data before fine-tuning pays? Fewer than about 20 labelled images per
class: probe only, since fine-tuning overfits. Tens: probe, then unfreeze the
last block at a small rate. Hundreds: unfreeze progressively with differential
rates. Thousands and up: fine-tune everything, and consider training from
scratch. Flowers102 gives ten per class, which is the first row --- and is why
the frozen probe does so much of the work today. These are orders of magnitude
rather than thresholds, and the measurement that settles it for your corpus is
the probe against the fine-tune on the same validation split.

\subsection{The backbone, and its preprocessing}

ResNet-18: 11{,}689{,}512 parameters, trained on ImageNet, producing 512
numbers before its classifier.

\begin{failbox}[title={Failure condition --- the weights come with their own preprocessing}]
\code{ResNet18\_Weights.DEFAULT.transforms()} specifies a resize to 256, a
centre crop to 224, and per-channel normalisation with mean
$[0.485, 0.456, 0.406]$ and standard deviation $[0.229, 0.224, 0.225]$.

Use \emph{these} statistics, not the ones you computed in the previous lecture.
The network's first layer expects inputs on the scale it was trained with, and
substituting your own is a silent, uncrashing degradation.
\end{failbox}

Freezing is two lines --- set \code{requires\_grad = False} on every parameter,
then replace \code{net.fc}, because a newly constructed module has
\code{requires\_grad=True} and so the head is trainable and nothing else is.
That leaves 52{,}326 trainable parameters.

What \code{requires\_grad = False} actually does: it stops the tensor being a
leaf that accumulates a gradient, so \code{backward()} does not fill
\code{p.grad}. The backward pass still travels through those layers if anything
downstream needs it --- here nothing does. And it is \emph{not} the same as
\code{eval()}.

\begin{failbox}[title={Failure condition --- freezing weights does not freeze a network}]
Batch-norm running statistics are \emph{buffers}, updated in the forward pass
whenever the module is in training mode. No gradient is involved, so
\code{requires\_grad = False} does nothing to them.

A ``frozen'' backbone in \code{train()} mode quietly replaces ImageNet's
statistics with the statistics of 1{,}020 flower photographs, in batches of 32.
The weights are fixed and the network still drifts --- towards a corpus of
1{,}020 images, which is the opposite of what you froze it for.

The fix is \code{eval()} on the frozen part: the same call as Lecture 10's
dropout bug, for the same underlying reason.
\end{failbox}

If nothing before the head is training, the backbone's output for a given image
never changes --- so run it once and fit the head on cached features. That is
not an optimisation; it is a consequence of freezing. Five seconds for the
features and eight for sixty epochs of the head.

\begin{center}
\small
\begin{tabular}{lrr}
\toprule
\textbf{Model} & \textbf{Test accuracy} & \textbf{Wall clock} \\
\midrule
From scratch, 80 epochs & 15.3\% & 364 s \\
Frozen backbone, linear head & 81.2\% & 13 s \\
\bottomrule
\end{tabular}
\end{center}

27 times faster and 66 points better. Why a linear model on somebody else's
features beats your network: the 52{,}326 trainable parameters are fitted with
1{,}020 examples, about 51 parameters per image against 4713 before, while the
11{,}176{,}512 frozen ones were fitted with 1.28 million examples by somebody
else.

\begin{keybox}[title={The variance problem was not solved}]
It was \emph{moved} --- to a dataset large enough to absorb it.
\end{keybox}

\subsection{Letting the last block move}

The early layers hold edges and colours, which are not about flowers; the late
layers hold parts and textures, which partly are. Unfreeze \code{layer4} and
give it a learning rate an order of magnitude below the head's. One optimiser,
two parameter groups --- that is all a differential learning rate is.

\begin{center}
\small
\begin{tabular}{lr}
\toprule
\textbf{Part of resnet18} & \textbf{Parameters} \\
\midrule
\code{conv1} through \code{layer3}, frozen & 2{,}782{,}784 \\
\code{layer4} and the new head, trained & 8{,}446{,}054 \\
\bottomrule
\end{tabular}
\end{center}

``Only the last block'' is not a small thing: three quarters of the parameters
are in the block we unfroze. That is Lecture 12's memory calculation mirrored
in a design decision --- parameters concentrate in the deep, low-resolution
layers, so freezing the \emph{general} layers freezes very few of them. What we
froze is the part that does general work, not the part that is large.

The two rates have to differ because the head starts random and has everything
to learn, while the pretrained block is already good and has something to lose.
Give the backbone the head's rate and the first few batches --- whose gradients
come through a randomly initialised head --- will damage the representation you
downloaded. It does not crash; the accuracy simply comes out lower and you
conclude that transfer learning does not work.

\begin{failbox}[title={Failure condition --- catastrophic forgetting}]
The pretrained weights being a good starting point is exactly what makes them
fragile: a step large enough to be useful on random weights is far too large
for weights that are already nearly right.

Which is why the usual recipe warms up the new head first while everything else
is frozen.
\end{failbox}

\subsection{Augmentation, and what each transform asserts}

We cannot buy labels; we can show the network a different view of the same
labelled image every epoch. A random resized crop teaches scale and position
tolerance; a horizontal flip says a flower is still that species mirrored. Each
is a transformation \emph{the label is known to survive}, and that is the
entire criterion --- a modelling decision, not a library call.

\begin{keybox}[title={What augmentation is not}]
It is not more information: the mutual information between the corpus and the
labels is unchanged. It is a \emph{constraint} --- it tells the network the
answer must be the same across a family of inputs, which puts it in the same
box as weight sharing and as ridge. Lecture 5's second cure, arriving for the
third time.

And a transformation the label does \emph{not} survive is a labelling error you
manufactured.
\end{keybox}

\begin{center}
\small
\begin{tabular}{llp{0.28\textwidth}}
\toprule
\textbf{Transform} & \textbf{What it asserts} & \textbf{True for flowers?} \\
\midrule
horizontal flip & mirror image is the same class & yes \\
random resized crop & a part identifies the whole & usually \\
colour jitter & hue is not diagnostic & dangerous --- colour is a flower
                                        feature \\
vertical flip & upside-down is the same class & often, but not for scenes \\
rotation & orientation carries no information & yes here \\
\bottomrule
\end{tabular}
\end{center}

Augmentation trades bias for variance: too little and the model memorises
1{,}020 images, too much and the training distribution stops resembling the
test one --- and the symptom of too much is both curves flat and low, which
Lecture 5 named as bias. Start with flip and crop, and add one transform at a
time, each a hypothesis about an invariance that can be tested like anything
else.

\section{A failure condition: one transform, two loaders}

It is easy to build one pipeline and hand it to both loaders. The training
loader wants random crops and flips; the evaluation loader must not have them.

\begin{failbox}[title={Failure condition --- an augmented evaluation set is not a measurement}]
It makes the score a \emph{random variable} rather than a function of the
weights. Score one fixed model --- the fine-tuned one --- on the clean
validation set and it gives 90.88\%; on the augmented set, ten times, the mean
is 90.02\% with a spread of 1.57 points.

The same weights, the same set, 1.57 points apart. Lecture 10's reviewer
question arriving again: a deterministic function of fixed weights and fixed
data does not wobble.
\end{failbox}

\begin{center}
\small
\begin{tabular}{lr}
\toprule
\textbf{Consequence} & \textbf{Size, measured} \\
\midrule
The reported number is pessimistic & 0.86 points low \\
The number is not reproducible & 1.57 points of noise \\
Model selection uses a noisy criterion & a different epoch chosen \\
\bottomrule
\end{tabular}
\end{center}

The clean validation curve peaks at epoch 14 and the augmented one at epoch 12,
so early stopping on the second keeps the wrong checkpoint. The first
consequence is harmless; the third is not, and it is the one nobody notices,
\emph{because a pessimistic number feels like the safe kind of wrong}.

\begin{keybox}[title={The check, in two lines}]
\begin{verbatim}
a = evaluate(model, val_loader)
b = evaluate(model, val_loader)
assert a == b, f"evaluation is not deterministic: {a} vs {b}"
\end{verbatim}
A comment saying ``do not augment the validation set'' is read once. This runs
every time, costs two forward passes, and fails loudly at the moment the
mistake is made rather than three lectures later --- and it catches all three
causes: augmentation, dropout still active, and a shuffled loader whose batches
are averaged rather than pooled.
\end{keybox}

Three more for the silent-failure catalogue, none of which raises: normalisation
statistics computed over the whole dataset, worth 0.25 points here and more on
a smaller corpus; augmentation applied to the validation set, worth 1.57 points
of noise; and a frozen backbone left in \code{train()} mode, whose statistics
drift with no symptom at all.

\section{The ablation}

\begin{center}
\small
\begin{tabular}{lrr}
\toprule
\textbf{Change} & \textbf{Test accuracy} & \textbf{Gain} \\
\midrule
from scratch & 15.3\% & --- \\
+ pretrained backbone, frozen & 81.2\% & 65.9 pts \\
+ \code{layer4} unfrozen, no augmentation & 86.3\% & 5.1 pts \\
+ augmentation & 88.0\% & 1.7 pts \\
\bottomrule
\end{tabular}
\end{center}

Almost all of the gain is in one row, and saying \emph{which} row is what makes
the result useful to somebody else.

\begin{failbox}[title={Failure condition --- what these numbers are not}]
All three transfer rows are single runs. Lecture 12 measured a seed-to-seed
spread of 2.34 points on this dataset, so the 5.1 is a difference and the 1.7
is not --- and neither has been replicated.

The wall clocks are for one laptop; the ratios transfer between machines much
better than the seconds do. And the pretrained weights saw 1.28 million
labelled photographs: that labelling cost is real, it is simply not ours.
\end{failbox}

\begin{center}
\small
\begin{tabular}{lrr}
\toprule
\textbf{Model} & \textbf{Test accuracy} & \textbf{Wall clock} \\
\midrule
Uniform guess & 0.98\% & --- \\
Commonest species & 3.87\% & --- \\
Convolutional net, from scratch & 15.3\% & 364 s \\
Frozen backbone, linear head & 81.2\% & 13 s \\
Fine-tuned + augmented & 88.0\% & 176 s \\
\midrule
Operator's requirement & 90\% & --- \\
\bottomrule
\end{tabular}
\end{center}

73 points of accuracy and 27 times less wall clock for the probe alone. Be
precise about ``cheaper'', though: the fine-tune is only 2.1 times faster than
from scratch, because it runs a much larger network on much larger images. The
frozen probe is the one that is dramatically cheaper --- and it is also the one
you should run first, every time.

\subsection{When the recipe does not work}

The domains genuinely differ --- medical scans, satellite bands, spectrograms
and line drawings are not photographs, and the middle blocks stop being useful
sooner. The input is not an image at all: reshaping tabular data into a square
does not make convolution appropriate, because there is no locality to exploit.
You have plenty of labels, in which case training from scratch catches up and
passes a frozen backbone --- transfer is strongest exactly where labels are
scarce. Or the task needs resolution the backbone discarded, which is Lecture
14's whole problem.

How to tell, cheaply: fit a linear head on frozen features and compare with the
trivial baseline. If the frozen features barely beat it, the backbone does not
understand your images and no amount of fine-tuning will hide it.

\begin{keybox}[title={One line to keep}]
Before you train a vision model from scratch, spend 13 seconds measuring what a
frozen pretrained backbone and a linear head already give you. On this
application that number was 81.2\% against 15.3\%.

You are entitled to train from scratch. You are not entitled to do it without
knowing what you turned down.
\end{keybox}

And what we still cannot do: say \emph{where} the flower is, because we have
one label per photograph and pooling threw the position away; handle two
flowers in one frame; or say ``I do not know'', which was half the operator's
original request.

\section{The notebook}

\begin{trybox}[title={What to change}]
Freeze everything including the final block and refit. The accuracy drops by a
few points --- which tells you how much of the gain came from adapting the last
block rather than from the backbone itself, and separates the two rows of the
ablation that are easiest to conflate.
\end{trybox}

\section{Exercises}

\begin{enumerate}[leftmargin=1.6em,itemsep=4pt]
  \item A network fits its training set to 71.8\% and its validation set to
        13.7\%. Name the condition, and say which of Lecture 5's two cures is
        available when labelling is the dominant cost. \hfill \scope{[4]}
  \item Explain why a linear head on frozen features can beat a network trained
        end to end on the same data, in terms of parameters per labelled
        example. \hfill \scope{[5]}
  \item \code{requires\_grad = False} is set on every parameter of a backbone,
        and something in it changes anyway during training. State what, why,
        and the fix. \hfill \scope{[5]}
  \item Give two augmentations that are safe for photographs of flowers and one
        that is not, with the claim about the task each one makes.
        \hfill \scope{[4]}
  \item An evaluation pipeline shares its transform with the training loader.
        Name the three consequences, and say which is the dangerous one and
        why. \hfill \scope{[4]}
\end{enumerate}

\section*{Summary}
\addcontentsline{toc}{section}{Summary}

Nothing was wrong. The loop, the metric, the split and the architecture were
all correct, and the network still reached 15.3\% --- because 1{,}020 labelled
images were being asked to pay for 4{,}807{,}494 parameters, most of which
encode facts about photographs rather than about flowers. The gap of 58 points
between training and validation is Lecture 5's variance, and with labelling as
the dominant cost the only available cure is more constraint.

A frozen backbone and a linear head reached 81.2\% in 13 seconds --- 66 points
better and 27 times cheaper --- because the trainable parameters fell to about
51 per labelled image while 11{,}176{,}512 frozen ones had been fitted on 1.28
million photographs by somebody else. The variance problem was not solved but
moved. Unfreezing the last block added 5.1 points and augmentation 1.7, against
a seed-to-seed spread of 2.34 that makes only the first of those a difference.

Three failures on the way, none of which raises: preprocessing with your own
normalisation instead of the weights' own; a frozen backbone left in
\code{train()} mode, whose batch-norm buffers drift towards 1{,}020 flowers
because they are not parameters and freezing gradients does not touch them; and
one transform handed to both loaders, which cost 1.57 points of noise and
selected a different epoch.

\end{document}
